%% Question bank chapter file - generated by bank_tools.py build (do not edit while the bank is still being built:
%% it is regenerated from the per-test data in Question Bank\data). Every question is one bankq block:
%%   \begin{bankq}[id=..., source={...}, paper=.., qno=.., type=..., subtopic={...}, answer={...}, answer_src=key|solved,
%%                 working={...}, difficulty=1-10, diff_parts={K S I C F}, diff_note={...},
%%                 group=<passage id or empty>, also_in={duplicate copies in other tests}]
%% difficulty = min(10, K+S+I+C+F-1): K concepts 1-3, S steps 1-4, I insight 0-2, C computation 0-2, F format 0-1
%%   ...question LaTeX, options, figures (\bfig / \bfigi -> img/)...
%%   \end{bankq}
%% Passages shared by several questions are bankpassage blocks; their questions carry group=<passage id>.
\documentclass[10pt]{article}
\usepackage[a4paper,margin=1.8cm]{geometry}
\usepackage[T1]{fontenc}
\usepackage{amsmath,amssymb,mathtools}
\usepackage{graphicx}
\usepackage[export]{adjustbox}
\usepackage{xcolor,needspace,keyval}
\usepackage{enumitem}
\usepackage{booktabs,array,tabularx}
\usepackage[hidelinks]{hyperref}
\setlength{\parindent}{0pt}
\setlength{\parskip}{3pt}
\makeatletter
\define@key{bq}{id}{\def\bq@id{#1}}\define@key{bq}{source}{\def\bq@src{#1}}\define@key{bq}{paper}{\def\bq@paper{#1}}
\define@key{bq}{qno}{\def\bq@qno{#1}}\define@key{bq}{type}{\def\bq@type{#1}}\define@key{bq}{subtopic}{\def\bq@sub{#1}}
\define@key{bq}{answer}{\def\bq@ans{#1}}\define@key{bq}{answer_src}{\def\bq@asrc{#1}}\define@key{bq}{working}{\def\bq@work{#1}}
\define@key{bq}{group}{\def\bq@group{#1}}\define@key{bq}{also_in}{\def\bq@also{#1}}
\define@key{bq}{difficulty}{\def\bq@diff{#1}}\define@key{bq}{diff_parts}{\def\bq@dparts{#1}}\define@key{bq}{diff_note}{\def\bq@dnote{#1}}
\newenvironment{bankq}[1][]{\def\bq@work{}\def\bq@also{}\def\bq@group{}\def\bq@diff{?}\def\bq@dparts{}\def\bq@dnote{}\setkeys{bq}{#1}%
  \par\needspace{4\baselineskip}\medskip\noindent{\small\textbf{\bq@id}\hfill\color{gray}[\bq@src, Paper \bq@paper, Q\bq@qno]
  \bq@type\ $\cdot$ \bq@sub\ $\cdot$ \textbf{D\bq@diff}}\par\nopagebreak\relax}
  {\par{\small\color{teal}Answer: \textbf{\bq@ans}\ifx\bq@work\@empty\else\ (solved: \bq@work)\fi
  \ifx\bq@also\@empty\else\par\color{gray}Also in: \bq@also\fi}\par}
\define@key{bp}{id}{\def\bp@id{#1}}\define@key{bp}{source}{\def\bp@src{#1}}\define@key{bp}{questions}{\def\bp@qs{#1}}
\newenvironment{bankpassage}[1][]{\setkeys{bp}{#1}\par\needspace{4\baselineskip}\bigskip\noindent
  {\small\textbf{Passage \bp@id}\hfill\color{gray}[\bp@src; questions \bp@qs]}\par\nopagebreak\relax\itshape}{\par}
\makeatother
\newcommand{\bfig}[2]{\par\smallskip{\centering\includegraphics[width=#1pt,max width=0.62\linewidth,max height=0.32\textheight]{img/#2}\par}\smallskip}
\newcommand{\bfigi}[2]{\includegraphics[width=#1pt,max width=\linewidth,max height=0.32\textheight,valign=c]{img/#2}}
\newcommand{\bmissing}[1]{\textit{[figure missing in the source conversion: \detokenize{#1}]}}
\newenvironment{qoptions}{\begin{itemize}[leftmargin=1.8em,itemsep=1pt,topsep=2pt,label={}]}{\end{itemize}}
\newcommand{\opt}[2]{\item[(#1)] #2}
\providecommand{\fwstart}{}\providecommand{\fwend}{}\providecommand{\qfullsplit}{}
\newcolumntype{L}[1]{>{\raggedright\arraybackslash}p{#1}}
\newcolumntype{C}[1]{>{\centering\arraybackslash}p{#1}}
\newcolumntype{R}[1]{>{\raggedleft\arraybackslash}p{#1}}
\providecommand{\sectionhdr}[1]{}\providecommand{\subsectionhdr}[1]{}


\begin{document}
{\Large\bfseries Aldehydes \& Ketones --- Question Bank}\par\smallskip
{\small Chemistry chapter C20. Questions from all converted test papers, grouped by source series. Answers: printed key where it matches the paper, otherwise solved (answer\_src=solved, with working). Exact repeats across papers appear once, with also\_in listing the other copies.}\par\medskip\hrule

\section*{Score-I papers}
\begin{bankq}[id=SI12-P1-CHEM-S1-Q01, source={Score-I - Test 12 (09 Mar 2022)}, paper=1, qno=1, type=multi, subtopic={Nucleophilic addition}, answer={A,C,D}, answer_src=key, difficulty=2, diff_parts={K1 S1 I0 C0 F1}, diff_note={concepts: nucleophilic addition reactivity | steps: compare the carbonyl compounds in options A and B; compare those in C and D}]
\bfig{127.6}{SI12_p1_chem_s1_q01_1.png}
\begin{center}\small
\begin{tabular}{@{}l@{\ }c@{\qquad\qquad}l@{\ }c@{}}
(A) & \bfigi{55.4}{SI12_p1_chem_s1_q01_A.png}\rule{0pt}{40pt} & (B) & CH$_3$--CHO\\[8pt]
(C) & \bfigi{58.2}{SI12_p1_chem_s1_q01_C.png}\rule{0pt}{26pt} & (D) & \bfigi{51.2}{SI12_p1_chem_s1_q01_D.png}
\end{tabular}
\end{center}
\end{bankq}
\begin{bankq}[id=SI12-P1-CHEM-S2-Q01, source={Score-I - Test 12 (09 Mar 2022)}, paper=1, qno=1, type=numerical, subtopic={Nucleophilic addition}, answer={4.00}, answer_src=key, difficulty=2, diff_parts={K1 S1 I0 C0 F1}, diff_note={concepts: iodoform test | steps: identify the methyl ketone or CH3CH(OH) groups; count the compounds}]
How many compound(s) gives Iodoform test with NaOH + I$_2$:
\bfig{311.5}{SI12_p1_chem_s2_q01_1.png}
\end{bankq}
\begin{bankq}[id=SI12-P2-CHEM-S2-Q02, source={Score-I - Test 12 (09 Mar 2022)}, paper=2, qno=2, type=numerical, subtopic={Nucleophilic addition}, answer={4.00}, answer_src=key, difficulty=2, diff_parts={K1 S1 I0 C0 F1}, diff_note={concepts: haloform reaction stoichiometry | steps: 3 equivalents for tribromination and 1 for cleavage = 4}]
How many equivalents of base are consumed in the following reaction ?\\
$\underset{\text{(1 equivalent)}}{\text{CH}_3\text{CHO}}\xrightarrow{\text{Br}_2/\bar{\text{O}}\text{H}/\Delta}\text{HCBr}_3+\text{HCOO}^-$
\end{bankq}

\section*{TOAS papers}
\begin{bankpassage}[id=TOAS04-P1-CHEM-S3-G5-6, source={TOAS - Test 04 (15 Mar 2022)}, questions={5-6}]
\bfig{89.9}{TOAS04_p1_chem_s3_para5_6_1.png}
\end{bankpassage}
\begin{bankq}[id=TOAS04-P1-CHEM-S3-Q05, source={TOAS - Test 04 (15 Mar 2022)}, paper=1, qno=5, type=numerical, subtopic={Nucleophilic addition}, answer={3.00}, answer_src=key, difficulty=2, diff_parts={K1 S1 I0 C0 F1}, diff_note={concepts: nucleophilic addition products | steps: identify the product; value of n}, group=TOAS04-P1-CHEM-S3-G5-6]
Value of `$n$' is?
\end{bankq}
\begin{bankq}[id=TOAS04-P1-CHEM-S3-Q06, source={TOAS - Test 04 (15 Mar 2022)}, paper=1, qno=6, type=numerical, subtopic={Nucleophilic addition}, answer={2.00}, answer_src=key, difficulty=2, diff_parts={K1 S1 I0 C0 F1}, diff_note={concepts: degree of unsaturation | steps: product formula; DOU}, group=TOAS04-P1-CHEM-S3-G5-6]
Degree of unsaturation (D.O.U.) in product ($X$)?
\end{bankq}
\begin{bankq}[id=TOAS07-P1-CHEM-S2-Q06, source={TOAS - Test 07 (01 May 2022)}, paper=1, qno=6, type=multi, subtopic={Nucleophilic addition}, answer={A,B,C}, answer_src=key, difficulty=3, diff_parts={K1 S2 I0 C0 F1}, diff_note={concepts: reactions of aldehydes with Grignard reagents and PCl5 | steps: options A and B; option C; option D}]
Out of the given options reaction(s) in which the substrate `X' can be acetaldehyde to give the final product :
\begin{qoptions}
\opt{A}{$\text{`X'}\xrightarrow[\text{(ii) NH}_4\text{Cl}]{\text{(i) CH}_3\text{MgBr}}\mathrm{CH_3{-}CH(OH){-}CH_3}$}
\opt{B}{$\text{`X'}\xrightarrow{\mathrm{PCl_5}}\mathrm{CH_3{-}CHCl_2}$}
\opt{C}{$\text{`X'}\xrightarrow[\text{(ii) HOH}]{\text{(i) Mg(Hg)}}B\xrightarrow[\Delta]{\text{Conc. H}_2\text{SO}_4}\mathrm{CH_3{-}CO{-}CH_2{-}CH_3}$}
\opt{D}{$\text{`X'}\xrightarrow[\text{(ii) Br}_2/\text{CCl}_4]{\text{(i) AgOH}}\mathrm{CH_3Br}$}
\end{qoptions}
\end{bankq}
\begin{bankq}[id=TOAS08-P1-CHEM-S3-Q06, source={TOAS - Test 08 (15 May 2022)}, paper=1, qno=6, type=numerical, subtopic={Nucleophilic addition}, answer={3}, answer_src=solved, working={Aldehydes/ketones C$_5$H$_{10}$O giving a chiral alcohol from an achiral carbonyl (racemic): pentanal, 3-methylbutanal, 2,2-dimethylpropanal; 2-methylbutanal gives diastereomers, ketones give achiral alcohols: 3}, difficulty=4, diff_parts={K1 S2 I1 C0 F1}, diff_note={concepts: carbonyl isomers and Grignard products | steps: isomers of C5H10O; which carbonyls give a chiral secondary alcohol from an achiral carbonyl; count | I: chirality created by addition to a prochiral carbonyl}]
If all carbonyls isomers of molecular mass $=86$ u are separately treated with $\mathrm{CH_3MgBr}$ followed by acid hydrolysis, how many of them will give racemic mixture ?
\end{bankq}
\begin{bankq}[id=TOAS08-P2-CHEM-S3-Q08, source={TOAS - Test 08 (15 May 2022)}, paper=2, qno=8, type=numerical, subtopic={Nucleophilic addition}, answer={3.00}, answer_src=key, difficulty=3, diff_parts={K1 S2 I0 C0 F1}, diff_note={concepts: reductions of carbonyl compounds | steps: identify the required conversion; check each reagent; count}]
How many of the following are suitable reagents for following conversion :
\bfig{98.2}{TOAS08_p2_chem_s3_q08_1.png}
(a) (i) $\mathrm{Ni/H_2}$ (ii) TsCl (iii) $\mathrm{NaBH_4/H_2O}$\newline
(b) $\mathrm{ZnHg/HCl}$\newline
(c) (i) $\mathrm{N_2H_4}$ / (ii) $\mathrm{\overset{\ominus}{O}H}$\newline
(d) $\mathrm{N_2H_4/H_2O_2}$\newline
(e) Na / Liq. $\mathrm{NH_3}$\newline
(f) $\mathrm{CH_3MgBr}$ followed by LAH\newline
(g) (i) DIBAL-H ($-78^\circ$C) , (ii) $\mathrm{H_2O}$
\end{bankq}
\begin{bankq}[id=TOAS09-P2-CHEM-S1-Q12, source={TOAS - Test 09 (30 May 2022)}, paper=2, qno=12, type=single, subtopic={Nucleophilic addition}, answer={A}, answer_src=key, difficulty=1, diff_parts={K1 S1 I0 C0 F0}, diff_note={concepts: product of a reaction | steps: identify the reaction; product}]
\bfig{109.2}{TOAS09_p2_chem_s1_q12_1.png}
Product P is?
\begin{qoptions}
\opt{A}{\bfigi{60.0}{TOAS09_p2_chem_s1_q12_A.png}}
\opt{B}{\bfigi{54.6}{TOAS09_p2_chem_s1_q12_B.png}}
\opt{C}{\bfigi{48.5}{TOAS09_p2_chem_s1_q12_C.png}}
\opt{D}{\bfigi{58.0}{TOAS09_p2_chem_s1_q12_D.png}}
\end{qoptions}
\end{bankq}
\begin{bankq}[id=TOAS10-P1-CHEM-S1-Q06, source={TOAS - Test 10 (12 Jun 2022)}, paper=1, qno=6, type=multi, subtopic={Nucleophilic addition}, answer={A}, answer_src=key, difficulty=3, diff_parts={K1 S2 I0 C0 F1}, diff_note={concepts: bisulphite addition compounds | steps: C3H6O giving the haloform test is acetone; bisulphite adduct structure; crystalline nature}]
(P) + NaHSO$_3\rightleftharpoons$ (Q)\par
Compound (P) having molecular formula `C$_3$H$_6$O' \& gives (+ve) haloform test. Choose the correct statement(s)?
\begin{qoptions}
\opt{A}{(Q) is \bfigi{81.3}{TOAS10_p1_chem_s1_q06_1.png} and crystalline in nature}
\opt{B}{(Q) is \bfigi{83.2}{TOAS10_p1_chem_s1_q06_2.png} and non-crystalline in nature}
\opt{C}{(P) can be aldehyde, methyl ketone \& some cyclic ketones as \bfigi{24.5}{TOAS10_p1_chem_s1_q06_4.png}}
\opt{D}{(Q) is \bfigi{88.6}{TOAS10_p1_chem_s1_q06_3.png} and crystalline in nature}
\end{qoptions}
\end{bankq}
\begin{bankq}[id=TOAS10-P1-CHEM-S1-Q08, source={TOAS - Test 10 (12 Jun 2022)}, paper=1, qno=8, type=multi, subtopic={Nucleophilic addition}, answer={B,D}, answer_src=key, difficulty=3, diff_parts={K1 S2 I0 C0 F1}, diff_note={concepts: haloform reaction sequence | steps: identify A and the iodoform product; products E and F; iodine consumption}]
\bfig{245.2}{TOAS10_p1_chem_s1_q08_1.png}
Correct statement is/are :
\begin{qoptions}
\opt{A}{Compound F is \bfigi{25.2}{TOAS10_p1_chem_s1_q08_2.png}}
\opt{B}{Compound E is \bfigi{22.5}{TOAS10_p1_chem_s1_q08_3.png}}
\opt{C}{A is acetophenone}
\opt{D}{In reaction from A to B, 3 equivalent of I$_2$ consumed, with respect to 1 equivalent of A.}
\end{qoptions}
\end{bankq}
\begin{bankq}[id=TOAS11-P1-CHEM-S3-Q03, source={TOAS - Test 11 (27 Jun 2022)}, paper=1, qno=3, type=numerical, subtopic={Nucleophilic addition}, answer={5.00}, answer_src=key, difficulty=4, diff_parts={K1 S2 I1 C0 F1}, diff_note={concepts: isomeric ketones and racemic products | steps: list the ketones of C6H12O with stereoisomers; reduction gives a chiral alcohol; count racemic products | I: chiral ketone gives diastereomers}]
Consider all possible isomeric ketones including stereoisomers of molecular formula C$_6$H$_{12}$O, All these isomers are independently reacted with NaBH$_4$ (\textbf{NOTE} : stereoisomers are also reacted separately). The total number of ketones that give a racemic product(s) is/are
\end{bankq}
\begin{bankq}[id=TOAS11-P2-CHEM-S3-Q04, source={TOAS - Test 11 (27 Jun 2022)}, paper=2, qno=4, type=single, subtopic={Nucleophilic addition}, answer={7.00}, answer_src=key, difficulty=2, diff_parts={K1 S2 I0 C0 F0}, diff_note={concepts: products of two reactions | steps: reaction A product P; reaction B product X; molar masses; difference}]
For the given reactions :\par
(A) \bfigi{174.4}{TOAS11_p2_chem_s3_q04_A.png}\par
(B) \bfigi{175.6}{TOAS11_p2_chem_s3_q04_B.png}\par
If the molar mass of P $=a$ and the molar mass of X $=b$\par
The value of $(a-b)/2=$ ?
\end{bankq}
\begin{bankq}[id=TOAS11-P2-CHEM-S3-Q05, source={TOAS - Test 11 (27 Jun 2022)}, paper=2, qno=5, type=single, subtopic={Nucleophilic addition}, answer={5.00}, answer_src=key, difficulty=1, diff_parts={K1 S1 I0 C0 F0}, diff_note={concepts: stable hydrates of carbonyl compounds | steps: electron withdrawing groups stabilise gem-diols; check each compound}]
How many of the following will form stable hydrate :\par
(1) CH$_3$COCH$_3$\quad(2) CCl$_3$CHO\quad(3) CF$_3$CHO\quad(4)
\bfig{270.8}{TOAS11_p2_chem_s3_q05_1.png}
\end{bankq}
\begin{bankq}[id=TOAS12-P1-CHEM-S1-Q05, source={TOAS - Test 12 (10 Jul 2022)}, paper=1, qno=5, type=multi, subtopic={Nucleophilic addition}, answer={A}, answer_src=key, difficulty=2, diff_parts={K1 S1 I0 C0 F1}, diff_note={concepts: products of a carbonyl reaction | steps: reactants and conditions; products A and B}]
\bfig{134.2}{TOAS12_p1_chem_s1_q05_A.png}
A and B are :-
\begin{qoptions}
\opt{A}{\bfigi{98.5}{TOAS12_p1_chem_s1_q05_A2.png}}
\opt{B}{\bfigi{94.2}{TOAS12_p1_chem_s1_q05_B.png}}
\opt{C}{\bfigi{93.7}{TOAS12_p1_chem_s1_q05_C.png}}
\opt{D}{No formation of A and B}
\end{qoptions}
\end{bankq}
\begin{bankq}[id=TOAS12-P1-CHEM-S1-Q08, source={TOAS - Test 12 (10 Jul 2022)}, paper=1, qno=8, type=multi, subtopic={Nucleophilic addition}, answer={A,B,C}, answer_src=key, difficulty=2, diff_parts={K1 S1 I0 C0 F1}, diff_note={concepts: selective product formation | steps: structure of the product in the figure; check the reactant in each option}]
Which of the following compound can only form
\bfig{164.1}{TOAS12_p1_chem_s1_q08_1.png}
\begin{qoptions}
\opt{A}{\bfigi{49.4}{TOAS12_p1_chem_s1_q08_A.png}}
\opt{B}{\bfigi{53.6}{TOAS12_p1_chem_s1_q08_B.png}}
\opt{C}{\bfigi{32.0}{TOAS12_p1_chem_s1_q08_C.png}}
\opt{D}{\bfigi{53.6}{TOAS12_p1_chem_s1_q08_D.png}}
\end{qoptions}
\end{bankq}
\begin{bankq}[id=TOAS12-P2-CHEM-S1-Q05, source={TOAS - Test 12 (10 Jul 2022)}, paper=2, qno=5, type=single, subtopic={Cannizzaro}, answer={B}, answer_src=key, difficulty=2, diff_parts={K1 S2 I0 C0 F0}, diff_note={concepts: reaction of a carbonyl compound with concentrated KOH | steps: absence of alpha hydrogen; Cannizzaro or crossed reaction; products}]
Which products are formed when compound \bfigi{50.0}{TOAS12_p2_chem_s1_q05_1.png} is treated with concentrated aqueous KOH solution?
\begin{qoptions}
\opt{A}{\bfigi{75.2}{TOAS12_p2_chem_s1_q05_A.png}}
\opt{B}{\bfigi{143.5}{TOAS12_p2_chem_s1_q05_B.png}}
\opt{C}{\bfigi{175.6}{TOAS12_p2_chem_s1_q05_C.png}}
\opt{D}{\bfigi{111.4}{TOAS12_p2_chem_s1_q05_D.png}}
\end{qoptions}
\end{bankq}
\begin{bankpassage}[id=TOAS12-P2-CHEM-S2-G9-10, source={TOAS - Test 12 (10 Jul 2022)}, questions={9-10}]
Aldehydes and ketones react with an alcohol from compounds called hemiacetals, in which there is one hydroxyl group and one either-like group. Reaction of a hemiacetal with a second molecule of alchol gives an acetal and a molecule of water. We study this reaction
\bfig{196.5}{TOAS12_p2_chem_s2_para9_10_1.png}
\noindent The carbonyl group of an aldehyde or ketone \quad A hemiacetal (has an --OH and an --OR group to the same carbon) \quad An acetal (has two --OR groups to the same carbon)\par
Draw structural formulas for the hemiacetal and acetal formed from these reagents. The stoichiometry of each reaction is given in the problem.\par
\noindent(A)
\bfig{139.6}{TOAS12_p2_chem_s2_para9_10_A.png}
\noindent(B)
\bfig{150.4}{TOAS12_p2_chem_s2_para9_10_B.png}
\noindent(C)
\bfig{150.4}{TOAS12_p2_chem_s2_para9_10_C.png}
\end{bankpassage}
\begin{bankq}[id=TOAS12-P2-CHEM-S2-Q09, source={TOAS - Test 12 (10 Jul 2022)}, paper=2, qno=9, type=multi, subtopic={Nucleophilic addition}, answer={A,C,D}, answer_src=key, difficulty=3, diff_parts={K1 S2 I0 C0 F1}, diff_note={concepts: hemiacetal and acetal formation | steps: reaction of the carbonyl with the alcohol; hemiacetal; acetal; which product is impossible}, group=TOAS12-P2-CHEM-S2-G9-10]
Product P is not possible -
\begin{qoptions}
\opt{A}{\bfigi{46.1}{TOAS12_p2_chem_s2_q09_A.png}}
\opt{B}{\bfigi{48.9}{TOAS12_p2_chem_s2_q09_B.png}}
\opt{C}{\bfigi{52.3}{TOAS12_p2_chem_s2_q09_C.png}}
\opt{D}{\bfigi{51.5}{TOAS12_p2_chem_s2_q09_D.png}}
\end{qoptions}
\end{bankq}
\begin{bankq}[id=TOAS12-P2-CHEM-S2-Q10, source={TOAS - Test 12 (10 Jul 2022)}, paper=2, qno=10, type=multi, subtopic={Nucleophilic addition}, answer={B,D}, answer_src=key, difficulty=2, diff_parts={K1 S1 I0 C0 F1}, diff_note={concepts: acetals of diols and carbonyl compounds | steps: products Q and R from the reagents; cyclic acetals}, group=TOAS12-P2-CHEM-S2-G9-10]
Which of the following statement is/are correct?
\begin{qoptions}
\opt{A}{Product Q is \bfigi{51.9}{TOAS12_p2_chem_s2_q10_1.png}}
\opt{B}{Product Q is \bfigi{65.6}{TOAS12_p2_chem_s2_q10_2.png}}
\opt{C}{Product R is \bfigi{41.9}{TOAS12_p2_chem_s2_q10_4.png}}
\opt{D}{Product R is \bfigi{46.6}{TOAS12_p2_chem_s2_q10_3.png}}
\end{qoptions}
\end{bankq}
\begin{bankq}[id=TOAS12-P2-CHEM-S3-Q03, source={TOAS - Test 12 (10 Jul 2022)}, paper=2, qno=3, type=numerical, subtopic={Cannizzaro}, answer={4.00}, answer_src=key, difficulty=3, diff_parts={K1 S2 I0 C0 F1}, diff_note={concepts: reactions correctly matched | steps: Cannizzaro with D2O; other reactions; count}]
Identify number of reactions correctly match with their products ?\par
(a) PhCH$=$O $\xrightarrow[\text{D}_2\text{O}]{\text{NaOD}}$ PhCH$_2$OD $+$ PhCO$_2^-$
\bfig{119.3}{TOAS12_p2_chem_s3_q03_1.png}
\bfig{237.5}{TOAS12_p2_chem_s3_q03_2.png}
\bfig{156.3}{TOAS12_p2_chem_s3_q03_3.png}
\end{bankq}
\begin{bankq}[id=TOAS13-P1-CHEM-S1-Q04, source={TOAS - Test 13 (24 Jul 2022)}, paper=1, qno=4, type=multi, subtopic={Cannizzaro}, answer={A,C,D}, answer_src=key, difficulty=5, diff_parts={K3 S2 I0 C0 F1}, diff_note={concepts: intramolecular Cannizzaro, benzilic acid rearrangement, Reformatsky | steps: glyoxal gives glycolate; benzil gives benzilic acid; Reformatsky with ester and acid work-up; check the fourth option}]
Select reaction(s) correctly matched with the major product :
\begin{qoptions}
\opt{A}{OHC--CHO $\xrightarrow[\Delta]{\text{Conc. NaOH}}$ HOCH$_2$--COO$^-$Na$^+$}
\opt{B}{\bfig{183.5}{TOAS13_p1_chem_s1_q04_B.png}\bfig{97.9}{TOAS13_p1_chem_s1_q04_3.png}}
\opt{C}{Ph--CO--CO--Ph $\xrightarrow[\text{(ii) H}^+]{\text{(i) OH}^-}$ Ph--C(OH)(Ph)--COOH}
\opt{D}{CH$_3$--CH=O $+$ Br--CH$_2$COOEt $\xrightarrow[\text{(ii) H}_3\text{O}^+\text{ (excess)}]{\text{(i) Zn / ether}}$ CH$_3$--CH(OH)--CH$_2$--COOH}
\end{qoptions}
\end{bankq}
\begin{bankq}[id=TOAS13-P1-CHEM-S3-Q04, source={TOAS - Test 13 (24 Jul 2022)}, paper=1, qno=4, type=numerical, subtopic={Aldol}, answer={6.00}, answer_src=key, difficulty=4, diff_parts={K1 S2 I1 C0 F1}, diff_note={concepts: crossed aldol enumeration | steps: enolisable partners are CH3CHO and the ketone; electrophiles are all three; two enolates times three electrophiles}]
The total number of all possible Mono Aldols ($\beta$-Hydroxy carbonyl compounds) are formed in the following reaction. (Ignore stereoisomers)\par
PhCHO $+$ CH$_3$CH=O $+$ Ph--CO--CH$_3$ $\xrightarrow{\text{dil. NaOH, }\Delta}$
\end{bankq}
\begin{bankq}[id=TOAS13-P2-CHEM-S1-Q10, source={TOAS - Test 13 (24 Jul 2022)}, paper=2, qno=10, type=single, subtopic={Cannizzaro}, answer={B}, answer_src=key, difficulty=2, diff_parts={K2 S1 I0 C0 F0}, diff_note={concepts: Rosenmund reduction, Cannizzaro reaction | steps: benzoyl chloride to benzaldehyde; concentrated NaOH disproportionates it to the alcohol and the carboxylate}]
The final products of the given reaction is :
\bfig{153.3}{TOAS13_p2_chem_s1_q10_1.png}
\begin{qoptions}
\opt{A}{\bfigi{133.9}{TOAS13_p2_chem_s1_q10_A.png}}
\opt{B}{\bfigi{135.4}{TOAS13_p2_chem_s1_q10_B.png}}
\opt{C}{\bfigi{134.9}{TOAS13_p2_chem_s1_q10_C.png}}
\opt{D}{\bfigi{114.6}{TOAS13_p2_chem_s1_q10_D.png}}
\end{qoptions}
\end{bankq}
\begin{bankq}[id=TOAS13-P2-CHEM-S1-Q11, source={TOAS - Test 13 (24 Jul 2022)}, paper=2, qno=11, type=single, subtopic={Aldol}, answer={A}, answer_src=key, difficulty=3, diff_parts={K1 S2 I1 C0 F0}, diff_note={concepts: intramolecular aldol condensation | steps: enolisation sites; the five-membered ring is favoured; dehydration to the enone}]
The major product obtained in the following reaction is :
\bfig{149.1}{TOAS13_p2_chem_s1_q11_1.png}
\begin{qoptions}
\opt{A}{\bfigi{49.9}{TOAS13_p2_chem_s1_q11_A.png}}
\opt{B}{\bfigi{60.0}{TOAS13_p2_chem_s1_q11_B.png}}
\opt{C}{\bfigi{52.4}{TOAS13_p2_chem_s1_q11_C.png}}
\opt{D}{\bfigi{58.3}{TOAS13_p2_chem_s1_q11_D.png}}
\end{qoptions}
\end{bankq}
\begin{bankq}[id=TOAS13-P2-CHEM-S2-Q04, source={TOAS - Test 13 (24 Jul 2022)}, paper=2, qno=4, type=numerical, subtopic={Cannizzaro}, answer={1}, answer_src=key, difficulty=5, diff_parts={K2 S2 I1 C0 F1}, diff_note={concepts: aldol with formaldehyde, crossed Cannizzaro | steps: three aldol steps to the trimethylol aldehyde; one crossed Cannizzaro gives pentaerythritol; count}]
The number of Cannizaro reaction(s) that occurs in the given transformation is
\bfig{99.9}{TOAS13_p2_chem_s2_q04_1.png}
\end{bankq}
\begin{bankq}[id=TOAS15-P1-CHEM-S1-Q01, source={TOAS - Test 15 (22 Aug 2022)}, paper=1, qno=1, type=single, subtopic={Aldol}, answer={B}, answer_src=key, difficulty=4, diff_parts={K2 S2 I1 C0 F0}, diff_note={concepts: retro-aldol and iodoform reaction, ketone derivative | steps: hydroxide cleaves the aldol to cyclopentanone and acetone; acetone gives iodoform with NaOI; the acid from the ketone and the final product}]
\bfig{161.3}{TOAS15_p1_chem_s1_q01_1.png}
in above reaction Y, A, B respectively is :-
\begin{qoptions}
\opt{A}{\bfigi{92.6}{TOAS15_p1_chem_s1_q01_A.png}}
\opt{B}{\bfigi{108.3}{TOAS15_p1_chem_s1_q01_B.png}}
\opt{C}{\bfigi{108.5}{TOAS15_p1_chem_s1_q01_C.png}}
\opt{D}{\bfigi{108.0}{TOAS15_p1_chem_s1_q01_D.png}}
\end{qoptions}
\end{bankq}
\begin{bankq}[id=TOAS15-P1-CHEM-S1-Q02, source={TOAS - Test 15 (22 Aug 2022)}, paper=1, qno=2, type=single, subtopic={Aldol}, answer={C}, answer_src=key, difficulty=3, diff_parts={K2 S2 I0 C0 F0}, diff_note={concepts: Michael addition, hydrolysis and decarboxylation | steps: enolate adds to the enone at the beta carbon; ester hydrolysis; decarboxylation of the malonic acid}]
\bfig{170.2}{TOAS15_p1_chem_s1_q02_1.png}
\begin{qoptions}
\opt{A}{\bfigi{74.2}{TOAS15_p1_chem_s1_q02_A.png}}
\opt{B}{\bfigi{74.2}{TOAS15_p1_chem_s1_q02_B.png}}
\opt{C}{\bfigi{79.4}{TOAS15_p1_chem_s1_q02_C.png}}
\opt{D}{\bfigi{68.9}{TOAS15_p1_chem_s1_q02_D.png}}
\end{qoptions}
\end{bankq}
\begin{bankq}[id=TOAS15-P1-CHEM-S2-Q11, source={TOAS - Test 15 (22 Aug 2022)}, paper=1, qno=11, type=multi, subtopic={Nucleophilic addition}, answer={A,C}, answer_src=key, difficulty=3, diff_parts={K2 S1 I0 C0 F1}, diff_note={concepts: Wittig reaction, ketone disconnection | steps: alkene disconnects to a ketone and an ylide; list the possible pairs}]
(P) \bfigi{33.5}{TOAS15_p1_chem_s2_q11_1.png}\bfigi{50.2}{TOAS15_p1_chem_s2_q11_2.png}\par
Possible P and Q is/are :-
\begin{qoptions}
\opt{A}{\bfigi{71.3}{TOAS15_p1_chem_s2_q11_A.png}}
\opt{B}{\bfigi{57.8}{TOAS15_p1_chem_s2_q11_B.png}}
\opt{C}{\bfigi{66.1}{TOAS15_p1_chem_s2_q11_C.png}}
\opt{D}{\bfigi{28.6}{TOAS15_p1_chem_s2_q11_D.png}\bfigi{28.9}{TOAS15_p1_chem_s2_q11_7D.png}}
\end{qoptions}
\end{bankq}
\begin{bankq}[id=TOAS16-P1-CHEM-S2-Q08, source={TOAS - Test 16 (04 Sep 2022)}, paper=1, qno=8, type=multi, subtopic={Nucleophilic addition}, answer={B,C,D}, answer_src=key, difficulty=6, diff_parts={K2 S3 I1 C0 F1}, diff_note={concepts: cyanohydrin formation, reduction and oxidation sequence | steps: HCN adds to the carbonyl; LAH reduces the nitrile; PCC oxidises; NaBH4 reduces; identify each product}]
Predict product (A), (B) (C) and (D)?
\bfig{209.4}{TOAS16_p1_chem_s2_q08_1.png}
\begin{qoptions}
\opt{A}{A is \bfigi{50.0}{TOAS16_p1_chem_s2_q08_2.png}}
\opt{B}{D is \bfigi{55.8}{TOAS16_p1_chem_s2_q08_3.png}}
\opt{C}{B is \bfigi{60.7}{TOAS16_p1_chem_s2_q08_4.png}}
\opt{D}{A is \bfigi{60.7}{TOAS16_p1_chem_s2_q08_5.png}}
\end{qoptions}
\end{bankq}
\begin{bankpassage}[id=TOAS16-P2-CHEM-S3-G5-6, source={TOAS - Test 16 (04 Sep 2022)}, questions={5-6}]
\bfig{240.4}{TOAS16_p2_chem_s3_para5_6_1.png}
\end{bankpassage}
\begin{bankq}[id=TOAS16-P2-CHEM-S3-Q05, source={TOAS - Test 16 (04 Sep 2022)}, paper=2, qno=5, type=numerical, subtopic={Nucleophilic addition}, answer={3.00}, answer_src=key, difficulty=5, diff_parts={K2 S2 I1 C0 F1}, diff_note={concepts: haloform reaction, resonance | steps: reagent for the haloform; anion of the by-product; count resonance structures}, group=TOAS16-P2-CHEM-S3-G5-6]
X react with hot and concentrated Y to give chloride and Z\par
Find total resonating structure of anion of Z.
\end{bankq}
\begin{bankq}[id=TOAS16-P2-CHEM-S3-Q06, source={TOAS - Test 16 (04 Sep 2022)}, paper=2, qno=6, type=numerical, subtopic={Nucleophilic addition}, answer={5.00}, answer_src=key, difficulty=5, diff_parts={K2 S2 I1 C0 F1}, diff_note={concepts: acetal formation, reduction | steps: acetal protection; LAH on the ketone-free acetal; identify D; principal chain length}, group=TOAS16-P2-CHEM-S3-G5-6]
How many carbon present in principle carbon chain of product (D) according to IUPAC rule?
\end{bankq}
\begin{bankq}[id=TOAS17-P2-CHEM-S1-Q08, source={TOAS - Test 17 (18 Sep 2022)}, paper=2, qno=8, type=single, subtopic={Aldol}, answer={D}, answer_src=key, difficulty=2, diff_parts={K1 S2 I0 C0 F0}, diff_note={concepts: named organic reactions | steps: aldol condensation; Reimer-Tiemann; check the others}]
Which of the following is incorrectly matched with their product?
\begin{qoptions}
\opt{A}{$2\,\mathrm{CH_3{-}CO{-}CH_3}\xrightarrow{\ominus\mathrm{OH}/\Delta}$ mesityl oxide}
\opt{B}{\bfigi{24.0}{TOAS17_p2_chem_s1_q08_2.png}\ $\xrightarrow{\mathrm{CHCl_3+NaOH}}$ Salicylaldehyde}
\opt{C}{\bfigi{128.3}{TOAS17_p2_chem_s1_q08_C.png}}
\opt{D}{\bfigi{83.5}{TOAS17_p2_chem_s1_q08_D.png}}
\end{qoptions}
\end{bankq}
\begin{bankq}[id=TOAS17-P2-CHEM-S1-Q11, source={TOAS - Test 17 (18 Sep 2022)}, paper=2, qno=11, type=single, subtopic={Cannizzaro}, answer={C}, answer_src=key, difficulty=3, diff_parts={K1 S2 I1 C0 F0}, diff_note={concepts: crossed aldol and Cannizzaro | steps: aldol with formaldehyde repeats; final crossed Cannizzaro; product mixture}]
\par\noindent\hfill$\mathrm{CH_3{-}CH_2{-}CH{=}O+\underset{\text{(excess)}}{HCHO}}\xrightarrow{\ \overline{\mathrm{O}}\mathrm{H}\ }$\hfill\null\par
Which of the following product mixture is formed as a final major product
\begin{qoptions}
\opt{A}{\bfigi{91.3}{TOAS17_p2_chem_s1_q11_A.png}}
\opt{B}{\bfigi{74.5}{TOAS17_p2_chem_s1_q11_B.png}}
\opt{C}{\bfigi{94.8}{TOAS17_p2_chem_s1_q11_C.png}}
\opt{D}{\bfigi{75.7}{TOAS17_p2_chem_s1_q11_D.png}}
\end{qoptions}
\end{bankq}
\begin{bankq}[id=TOAS17-P2-CHEM-S2-Q01, source={TOAS - Test 17 (18 Sep 2022)}, paper=2, qno=1, type=single, subtopic={Nucleophilic addition}, answer={6}, answer_src=key, difficulty=3, diff_parts={K2 S2 I0 C0 F0}, diff_note={concepts: named reactions, products | steps: apply each reaction; compare with the given product; count}]
How many reaction correctly match with their major product :
\bfig{152.0}{TOAS17_p2_chem_s2_q01_1.png}
\end{bankq}
\end{document}
